% Source: https://artofproblemsolving.com/community/c3h2697016_amgm_two_variables_but_extended
% Stable identity: aops-topic-2697016-q01
Let $Q=\sqrt{\dfrac{a^2+b^2}{2}}$ and $A=\dfrac{a+b}{2}$.
\begin{itemize}
\item $Q^2-A^2=\dfrac{(a-b)^2}{4}\ge0$, so $Q\ge A$. Then $\dfrac{a^2+b^2}{a+b}=\dfrac{Q^2}{A}\ge Q$.
\item Logarithmic mean versus $A$ and $\sqrt{ab}$: by symmetry let $b>a$ and $t=\dfrac ba>1$; the three means are $a\cdot\dfrac{1+t}{2}$, $a\cdot\dfrac{t-1}{\ln t}$ and $a\sqrt t$.
\begin{itemize}
\item $\dfrac{1+t}{2}\ge\dfrac{t-1}{\ln t}$ means $h(t)=\ln t-\dfrac{2(t-1)}{t+1}\ge0$. Here $h(1)=0$ and $h'(t)=\dfrac1t-\dfrac{4}{(t+1)^2}=\dfrac{(t-1)^2}{t(t+1)^2}\ge0$.
\item With $x=\sqrt t>1$, $\dfrac{t-1}{\ln t}\ge\sqrt t$ means $g(x)=\dfrac12\left(x-\dfrac1x\right)-\ln x\ge0$. Here $g(1)=0$ and $g'(x)=\dfrac{(x-1)^2}{2x^2}\ge0$.
\end{itemize}
\item $\sqrt{ab}\ge\dfrac{2ab}{a+b}$ is equivalent to $a+b\ge2\sqrt{ab}$, i.e. AM--GM.
\end{itemize}
$\blacksquare$ (All inequalities are strict for $a\ne b$; for $a=b$ all the means equal $a$, the logarithmic mean being defined as $a$.)
